\chapter{Critical-Line Phase Normalization and the Hardy Z Function}
\label{ch:hardy}

\section{Pure-phase duality on the fixed axis}
Write the functional equation as
\begin{equation}
\zeta(s)=\chi(s)\zeta(1-s).
\end{equation}
On $s=1/2+it$, one has $1-s=\overline{s}$, hence
\begin{equation}
\zeta\!\left(\frac12+it\right)
=\chi\!\left(\frac12+it\right)
\overline{\zeta\!\left(\frac12+it\right)}.
\label{eq:phase-duality}
\end{equation}
The modulus of $\chi$ is one on this line. Therefore there exists a real phase $\vartheta(t)$ such that
\begin{equation}
\chi\!\left(\frac12+it\right)=e^{-2i\vartheta(t)}.
\end{equation}
Define the Hardy function
\begin{equation}
Z(t)=e^{i\vartheta(t)}\zeta\!\left(\frac12+it\right).
\label{eq:hardyZ}
\end{equation}
Equation \eqref{eq:phase-duality} implies $Z(t)=\overline{Z(t)}$, so $Z(t)$ is real for real $t$. Its real zeros are precisely the zeros of zeta on the critical line \cite{Titchmarsh1986,DLMF25}.

This is an established phase-normalization result. It shows that the fixed axis converts the functional-equation duality into a real observable after removing a known phase.

\section{Balance of modulus and phase}
The critical line is distinguished by
\begin{equation}
|\chi(1/2+it)|=1.
\end{equation}
Thus duality changes phase but not magnitude. Away from the line, one generally has $|\chi(s)|\neq1$. In a two-branch interpretation, this means that the functional equation relates branches of unequal norm away from the fixed axis, while the axis supports unitary phase matching.

This statement must be handled carefully. A complex equation can vanish through cancellation even when individual terms have unequal raw norms if other coefficients or remainders are present. The modulus-one property is therefore a strong motivation for a balanced-state model, not an independent proof of RH.

\section{Approximate functional equation}
A standard approximate functional equation decomposes zeta into a direct Dirichlet sum and a dual sum multiplied by $\chi(s)$, schematically
\begin{equation}
\zeta(s)
=\sum_{n\leq x}n^{-s}
+\chi(s)\sum_{n\leq y}n^{s-1}
+R(s;x,y),
\qquad xy\asymp\frac{|t|}{2\pi}.
\label{eq:afe}
\end{equation}
On the critical line and with symmetric truncation $x\approx y$, the two sums are conjugate after phase normalization. This is the analytic-number-theory analogue of two complementary branches with matched weight.

The Riemann--Siegel formula sharpens this decomposition and is the basis of efficient evaluation of $Z(t)$. The observed zeros occur when the phase-normalized real quantity changes sign or reaches zero. The information-spinor programme proposes that the exact completion of this dual-sum structure may possess a projective two-channel interpretation.

\section{Riemann Xi as a real Fourier transform}
The completed function on the critical line admits a cosine-transform representation of the form
\begin{equation}
\Xi(t)=\int_0^\infty \Phi(u)\cos(tu)\,\dd u,
\label{eq:xi-fourier}
\end{equation}
for a rapidly decaying real kernel $\Phi$. Hence $\Xi$ is real and even. The Riemann hypothesis is equivalent to all zeros of this entire function occurring at real $t$.

This representation suggests a second route to the missing theorem. Instead of building a two-dimensional state for each zero, one may seek a positivity or total-positivity property of $\Phi$ strong enough to enforce real zeros of its Fourier transform. Classical work surrounding the de Bruijn--Newman deformation illustrates how heat-flow and zero geometry can encode the problem \cite{deBruijn1950,Newman1976,RodgersTao2020}.

\section{Phase quantization at a zero}
At a simple critical-line zero $t_n$, the real function $Z(t)$ crosses zero. Between zeros it has a sign, and its phase-normalized zeta value is either aligned or anti-aligned with the real readout. The cancellation theorem suggests representing the zero crossing by a relative phase of $\pi$ between dual contributions.

A rigorous implementation would require an exact identity
\begin{equation}
Z(t)=\mathcal N(t)\left[
\sqrt{\sigma(t)}F_+(t)
+e^{i\phi(t)}\sqrt{1-\sigma(t)}F_-(t)
\right]
\end{equation}
with $F_-$ canonically conjugate or complementary to $F_+$ and with nonvanishing normalization. On the critical line $\sigma=1/2$ would be fixed, and zerohood would reduce to phase opposition. The approximate functional equation provides a template but not yet such an exact normalized identity.

\begin{figure}[ht]
\centering
\includegraphics[width=0.84\textwidth]{hardy_z_zeros.pdf}
\caption{Numerical Hardy $Z(t)$ and the first ten tabulated critical-line zeros. This reproduces known data; it is not evidence that no off-line zeros exist.}
\label{fig:hardy}
\end{figure}
